% SPDX-License-Identifier: MIT
% Copyright (c) 2026 Evgeny Baulin
% Optimization Methods, Seminar 09: theory handout (English).
% Generated by checks/09. Linear Programming as the Foundation of Integer Optimization.ipynb from web/09/js/content-*.js and content/students.json. Do not edit by hand.
% Compile with: latexmk -xelatex -interaction=nonstopmode -halt-on-error Theory_en.tex
\documentclass[11pt,a4paper]{article}

\usepackage{fontspec}
\setmainfont{cmunrm.otf}[BoldFont=cmunbx.otf, ItalicFont=cmunti.otf, BoldItalicFont=cmunbi.otf]
\setsansfont{cmunss.otf}[BoldFont=cmunsx.otf, ItalicFont=cmunsi.otf, BoldItalicFont=cmunso.otf]
\setmonofont{cmuntt.otf}[BoldFont=cmuntb.otf, ItalicFont=cmunit.otf, BoldItalicFont=cmuntx.otf]
\usepackage{polyglossia}
\setdefaultlanguage{english}

\usepackage[a4paper,left=2.3cm,right=2.3cm,top=2.4cm,bottom=2.4cm,headheight=14pt]{geometry}
\usepackage{amsmath,amssymb}
\usepackage{graphicx}
\usepackage{array,booktabs,longtable}
\usepackage{enumitem}
\usepackage{xcolor}
\usepackage[most]{tcolorbox}
\usepackage{titlesec}
\usepackage{fancyhdr}
\usepackage{needspace}
\usepackage[expansion=false]{microtype}
\usepackage[colorlinks=true,linkcolor=hseblue,urlcolor=hseblue,citecolor=hseblue]{hyperref}

\definecolor{hseblue}{RGB}{16,45,105}
\definecolor{hseaccent}{RGB}{198,40,40}
\definecolor{hsegray}{RGB}{108,114,126}
\definecolor{hselight}{RGB}{240,243,248}
\definecolor{hsewarm}{RGB}{253,244,232}

\hypersetup{pdftitle={Optimization Methods. Seminar 09. Theory handout}, pdfauthor={Evgeny Baulin}}

\titleformat{\section}{\color{hseblue}\Large\bfseries}{\thesection}{0.6em}{}
\titleformat{\subsection}{\color{hseblue}\large\bfseries}{\thesubsection}{0.6em}{}
\titlespacing*{\section}{0pt}{16pt plus 4pt minus 2pt}{8pt plus 2pt}
\titlespacing*{\subsection}{0pt}{12pt plus 3pt minus 2pt}{6pt plus 2pt}
\pretocmd{\section}{\Needspace{7\baselineskip}}{}{\errmessage{cannot patch section}}
\pretocmd{\subsection}{\Needspace{5\baselineskip}}{}{\errmessage{cannot patch subsection}}

\tcbset{handoutbox/.style={enhanced, breakable, sharp corners, boxrule=0pt, left=7pt, right=7pt, top=5pt, bottom=5pt, before skip=8pt plus 2pt, after skip=8pt plus 2pt}}
\newtcolorbox{examplebox}{handoutbox, colback=hselight, colframe=hselight, before upper={\textcolor{hseblue}{\bfseries Example.}\ }}
\newtcolorbox{whybox}{handoutbox, colback=white, colframe=hseblue, leftrule=2.5pt, top=3pt, bottom=3pt, before upper={\textcolor{hseaccent}{\bfseries Why.}\ }}
\newtcolorbox{answerbox}{handoutbox, colback=hsewarm, colframe=hsewarm, before upper={\textcolor{hseaccent}{\bfseries Answer.}\ }}
\newtcolorbox{conventionbox}{handoutbox, colback=hselight, colframe=hseblue, leftrule=2.5pt, before upper={\textcolor{hseblue}{\bfseries Notation.}\par\smallskip}}

\newcommand{\handoutfigure}[3]{\par\medskip\noindent\begin{minipage}{\linewidth}\centering
  \includegraphics[width=#1\linewidth]{figures/#2.pdf}\par\smallskip
  \begin{minipage}{0.94\linewidth}\small\textcolor{hseblue}{\bfseries Figure.}\ #3\end{minipage}\end{minipage}\par\medskip}
\newcommand{\problemhead}[3]{\par\bigskip\noindent{\color{hseblue}\large\bfseries #1.\ #2}\par\nopagebreak
  \noindent{\small\itshape\color{hsegray}#3}\par\nopagebreak\smallskip}
\newtcolorbox{definitionbox}[1]{handoutbox, colback=white, colframe=hseblue, leftrule=2.5pt, before upper={\textcolor{hseblue}{\bfseries #1}\ }}
\newtcolorbox{theorembox}[1]{handoutbox, colback=hselight, colframe=hseblue, leftrule=2.5pt, before upper={\textcolor{hseblue}{\bfseries #1}\ }}
\newtcolorbox{proofbox}{handoutbox, colback=white, colframe=hsegray, leftrule=1.2pt, top=3pt, bottom=3pt, before upper={\textcolor{hsegray}{\bfseries\itshape Proof.}\ }}
\newtcolorbox{notebox}[1]{handoutbox, colback=hselight, colframe=hselight, before upper={\textcolor{hseblue}{\bfseries #1}\par\smallskip}}
\newtcolorbox{trapbox}{handoutbox, colback=hsewarm, colframe=hseaccent, leftrule=2.5pt, before upper={\textcolor{hseaccent}{\bfseries Trap.}\ }}
\newtcolorbox{selfcheckbox}{handoutbox, colback=white, colframe=hseaccent, leftrule=2.5pt, top=3pt, bottom=3pt, before upper={\textcolor{hseaccent}{\bfseries Check yourself.}\ }}
\newcommand{\handoutdisplay}[1]{\sbox0{$\displaystyle #1$}\ifdim\wd0>\linewidth
  \par{\centering\binoppenalty=10000\relpenalty=0\relax$\displaystyle #1$\par}\else\[#1\]\fi}
\newcommand{\selfcheckanswer}{\par\smallskip\textcolor{hsegray}{\bfseries Answer.}\ }

\setlist[itemize]{itemsep=3pt,topsep=4pt}
\setlist[enumerate]{itemsep=3pt,topsep=4pt}
\setlength{\parindent}{0pt}
\setlength{\parskip}{5pt plus 1pt}
\setlength{\emergencystretch}{2em}
\setcounter{tocdepth}{2}

\pagestyle{fancy}
\fancyhf{}
\fancyhead[L]{\small\color{hsegray}Optimization Methods}
\fancyhead[R]{\small\color{hsegray}Seminar 09}
\fancyfoot[L]{\small\color{hsegray}Evgeny Baulin}
\fancyfoot[C]{\small\thepage}
\fancyfoot[R]{\small\color{hsegray}MIT License}
\renewcommand{\headrulewidth}{0pt}
\renewcommand{\footrulewidth}{0pt}
\fancypagestyle{plain}{\pagestyle{fancy}}

\begin{document}

\begin{center}
  {\large\color{hsegray}Optimization Methods\par}
  \vspace{4pt}
  {\large\bfseries\color{hseaccent}Seminar 09\par}
  \vspace{8pt}
  {\LARGE\bfseries\color{hseblue}Linear Programming: Corners, the Simplex Method, Duality\par}
  \vspace{8pt}
  {\Large Theory handout\par}
  \vspace{12pt}
  HSE University, Faculty of Computer Science\par
  Academic year 2026-2027\par
  \vspace{8pt}
  Course lecturer: Andrey Ignatov\par
  Evgeny Baulin\par
  Telegram: \href{https://t.me/tarakan_tuc}{@tarakan\_tuc}\par
\end{center}

\noindent{\color{hseblue}\rule{\linewidth}{0.6pt}}

\tableofcontents

\clearpage


This handout follows seminar 9. It starts from the geometry of a linear program, turns corners into basic solutions, builds the simplex method on top of that, and then reads the dual solution as a list of prices. The last sections show how the method can fail and what all this gives to integer optimization. All numbers come from the checks notebook of the seminar.

\begin{conventionbox}
Primal: maximize $c^{\top}x$ subject to $Ax\le b$, $x\ge0$, with $A\in\mathbb R^{m\times n}$; $s=b-Ax$ are the slacks. Dual: minimize $b^{\top}y$ subject to $A^{\top}y\ge c$, $y\ge0$. In a tableau the last row is $z-c^{\top}x=0$, so a negative entry marks a variable worth raising. $B$ is a basis, $x_B$ the basic variables, $p^*(b)$ the optimal value as a function of the right-hand side.
\end{conventionbox}

% =====================================================================
\section{Linear programs and their corners}\label{sec:geometry}

A linear program optimizes a linear function over a set cut out by finitely many linear inequalities and equations. Such a set is a polyhedron: convex, with flat faces, possibly unbounded, possibly empty. An LP has exactly three possible outcomes: it is infeasible, it is unbounded, or it has an optimal solution, which is either unique or fills a whole face, in the plane an edge.

\begin{definitionbox}{Definition: Standard form.}
A problem $\min c^{\top}x$ subject to $Ax=b$, $x\ge0$. Any LP can be brought to it: add a slack to a row $\le$, subtract a surplus from a row $\ge$, write a free variable as $x^+-x^-$ with both parts nonnegative, multiply a row with $b_i<0$ by $-1$, and replace a maximization of $c^{\top}x$ by a minimization of $-c^{\top}x$.
\end{definitionbox}

\begin{definitionbox}{Definition: Basic solution.}
Let $A$ have $m$ independent rows. Choose $m$ independent columns, the basis $B$, set the remaining variables to zero and solve $Bx_B=b$. The result is a basic solution. It is a basic feasible solution (BFS) when $x_B\ge0$, and degenerate when some component of $x_B$ is zero.
\end{definitionbox}

\begin{theorembox}{Theorem: Corners are basic feasible solutions.}
A point of $\{x\ge0:Ax=b\}$ is a corner, that is, not a midpoint of two other feasible points, exactly when it is a basic feasible solution. If an LP in standard form has an optimum, it has an optimal corner.
\end{theorembox}

\begin{proofbox}
If $x$ is basic and $x=\tfrac{1}{2}(u+v)$ with feasible $u,v$, then $u$ and $v$ vanish wherever $x$ does, and $Bu_B=Bv_B=b$ gives $u=v=x$. Conversely, if the columns of the positive part of $x$ are dependent, a nonzero $d$ on that support has $Ad=0$, and $x\pm\varepsilon d$ are feasible, so $x$ is no corner. For the optimum: if an optimal $x$ has a dependent support, then $c^{\top}d=0$, since otherwise one of $x\pm\varepsilon d$ would beat $x$. Choose the sign of $d$ that gives it a negative component and move along it until a component hits zero. The value stays, the support shrinks, and after finitely many moves it is independent.
\end{proofbox}

\begin{whybox}
The level sets of a linear function are flat. Nothing pulls the optimum into the interior, so the last point of contact with the feasible set always contains a corner, and the search can be limited to finitely many points.
\end{whybox}

\handoutfigure{0.62}{fig_a1}{The blend of A1. Full dots are basic feasible solutions, crosses are basic solutions outside the feasible set.}

\problemhead{A1}{A feed blend}{On the board}
\textbf{Statement.} A feed is mixed from $x$ kg of a first ingredient at $5$ per kg and $y$ kg of a second at $4$ per kg. A batch needs at least $8$ units of protein and $9$ units of fibre; a kg of the first ingredient gives $2$ and $1$ of them, a kg of the second $1$ and $3$. Minimize the cost graphically, then bring the problem to standard form and list all its basic solutions. Give $x^*$, the cost $z^*$ and the number of feasible basic solutions.

\textbf{Solution.}
\begingroup\sloppy
\begin{enumerate}[label=\textbf{\arabic*.}, leftmargin=*]
  \item The model is $\min 5x+4y$ subject to $2x+y\ge8$, $x+3y\ge9$, $x,y\ge0$. The two lines meet at $(3,\ 2)$, and the feasible set is the unbounded region above both, with corners $(0,\ 8)$, $(3,\ 2)$ and $(9,\ 0)$.
  \item The level lines $5x+4y=z$ all have the same slope. Push one down towards the origin until it is about to lose the feasible set: the last point it touches is a corner. The costs at the corners are $(32,\ 23,\ 45)$, so the best is $(3,\ 2)$ with $z^*=23$.
  \item Standard form takes one surplus variable per row: $2x+y-s_1=8$, $x+3y-s_2=9$ with all four variables nonnegative. Two equations and four unknowns: pick two basic variables, set the other two to zero and solve. There are $\binom{4}{2}=6$ choices.
  \item All six at a glance:
\handoutdisplay{
\begin{array}{c|cccc|c|c} B & x & y & s_1 & s_2 & \ge 0 & z \\ \hline x,y & 3 & 2 & 0 & 0 & \checkmark & 23 \\ x,s_1 & 9 & 0 & 10 & 0 & \checkmark & 45 \\ x,s_2 & 4 & 0 & 0 & -5 & \times &  \\ y,s_1 & 0 & 3 & -5 & 0 & \times &  \\ y,s_2 & 0 & 8 & 0 & 15 & \checkmark & 32 \\ s_1,s_2 & 0 & 0 & -8 & -9 & \times &  \end{array}
}
 Three of the six have every variable nonnegative, and these are exactly the three corners. The other three are crossings of boundary lines outside the feasible set; a negative $s_i$ says which requirement they break.
\end{enumerate}
\endgroup
\begin{answerbox}
$(x^*,y^*)=(3,\ 2)$, $z^*=23$, 3 feasible basic solutions out of 6.
\end{answerbox}

\begin{trapbox}
The word basic refers to a choice of columns, not to a point. A degenerate corner belongs to several bases, and a basic solution may be infeasible. Counting bases does not count corners.
\end{trapbox}

\problemhead{A2}{All the basic solutions}{On your own}
\textbf{Statement.} Maximize $3x+2y$ subject to $x+y\le7$, $x+2y\le10$, $x\le5$ and $x,y\ge0$. Add three slack variables and go through every choice of three columns. Give the number of choices that yield a basic solution, the number of feasible ones and the optimal value.

\textbf{Solution.}
\begingroup\sloppy
\begin{enumerate}[label=\textbf{\arabic*.}, leftmargin=*]
  \item With slacks the system is $x+y+s_1=7$, $x+2y+s_2=10$, $x+s_3=5$. Choosing three of five columns can be done in $10$ ways. The choice $\{y,s_1,s_2\}$ leaves the third row without any basic variable: its columns are dependent. So the system has $9$ basic solutions.
  \item Solving each one, $5$ are nonnegative. They are the corners $(0,\ 0)$, $(5,\ 0)$, $(5,\ 2)$, $(4,\ 3)$ and $(0,\ 5)$ of the polygon; the rest have a negative slack.
  \item The objective at these corners takes the values $(0,\ 15,\ 19,\ 18,\ 10)$. The largest is $z^*=19$ at $(5,\ 2)$.
\end{enumerate}
\endgroup
\begin{answerbox}
9 basic solutions, 5 of them feasible, $z^*=19$ at $(5,\ 2)$.
\end{answerbox}

\problemhead{A3}{A free variable}{On your own}
\textbf{Statement.} Minimize $x+2y$ subject to $x+y\ge1$, $x\le6$ and $x\ge0$, with $y$ free in sign. Write the problem in standard form, replacing $y$ by a difference of two nonnegative variables. Give the number of variables of the standard form, then $y^*$ and $z^*$.

\textbf{Solution.}
\begingroup\sloppy
\begin{enumerate}[label=\textbf{\arabic*.}, leftmargin=*]
  \item With $y=y^+-y^-$ the system becomes $x+y^+-y^--s_1=1$, $x+s_2=6$: $5$ nonnegative variables $x,y^+,y^-,s_1,s_2$ and two rows. The objective is $x+2y^+-2y^-$.
  \item The objective falls as $y$ falls, so the first row is tight: $y=1-x$, and $z=x+2(1-x)=2-x$. That falls with $x$, so $x$ goes to its upper bound.
  \item So $x^*=6$, $y^*=-5$, $z^*=-4$. In standard form $y^+=0$ and $y^-=5$: the columns of $y^+$ and $y^-$ are opposite, so they never stand in one basis.
\end{enumerate}
\endgroup
\begin{answerbox}
5 variables, $y^*=-5$, $z^*=-4$.
\end{answerbox}

\begin{selfcheckbox}
Why can the columns of $y^+$ and $y^-$ of A3 never be basic together?
\selfcheckanswer They are opposite vectors, so they are dependent and cannot both belong to a basis.
\end{selfcheckbox}

% =====================================================================
\section{The simplex method}\label{sec:simplex}

Enumeration of bases is hopeless: the number $\binom{n}{m}$ explodes. The simplex method looks at very few of them. It starts at a corner and moves to a neighbouring corner with a better value until no neighbour is better. For a linear objective on a convex set a corner with no better neighbour is a global optimum.

\begin{definitionbox}{Definition: Tableau.}
The system $Ax+s=b$ together with $z-c^{\top}x=0$, rewritten so that each basic variable stands alone in its row. Each basic variable has a unit column, the last column shows their values, and the last row writes $z$ through the nonbasic variables. The entries $r_j$ of the $z$-row are the reduced costs: $z=z_0-\sum_jr_jx_j$ over the nonbasic $j$.
\end{definitionbox}

\begin{itemize}
  \item Entering: pick a nonbasic column with $r_j<0$, the most negative by Dantzig's rule. If there is none, stop: the point is optimal.
  \item Ratio test: among the rows with a positive entry in that column find the smallest ratio of the right-hand side to the entry. If no entry is positive, stop: the LP is unbounded.
  \item Pivot: divide that row by the entry and subtract multiples of it from the other rows, the $z$-row included, so that the column becomes a unit column.
\end{itemize}

\begin{theorembox}{Proposition: Optimality criterion.}
If every entry $r_j$ of the $z$-row is nonnegative at a basic feasible solution, that solution is optimal.
\end{theorembox}

\begin{proofbox}
The row $z=z_0-\sum_jr_jx_j$ is an identity on all solutions of the system. Any feasible point has $x_j\ge0$, so its value is at most $z_0$, and the current point reaches $z_0$. If some nonbasic $r_j=0$, pivoting $x_j$ in leaves $z$ as it is, and unless the step has length zero the edge between the two corners is optimal.
\end{proofbox}

\begin{whybox}
The ratio test keeps the point feasible: raising the entering variable to $t$ lowers each basic variable with a positive entry $a$ by $at$, and the first to reach zero stops the walk. That variable leaves the basis and a neighbouring corner is reached.
\end{whybox}

\handoutfigure{0.62}{fig_b1}{The path of B1: two pivots from the origin, along the $y$ axis and then along the top edge.}

\problemhead{B1}{A bakery by tableau}{On the board}
\textbf{Statement.} A bakery makes $x$ trays of bread and $y$ trays of pies; a tray earns $2$ and $5$. Oven time gives $x+y\le8$, dough gives $2x+y\le14$, and filling, needed for pies only, gives $2y\le10$. Solve by the simplex method with Dantzig's rule: the variable with the most negative entry of the $z$-row enters. Give $x^*$, $y^*$ and $z^*$.

\textbf{Solution.}
\begingroup\sloppy
\begin{enumerate}[label=\textbf{\arabic*.}, leftmargin=*]
  \item Add slacks $s_1,s_2,s_3$ and write $z-2x-5y=0$ as the last row. The basis is the slacks, the point is the origin, $z=0$. The boxed entry is the first pivot.
\handoutdisplay{
\begin{array}{c|ccccc|c}  & x & y & s_1 & s_2 & s_3 & \\ \hline s_1 & 1 & 1 & 1 & 0 & 0 & 8 \\ s_2 & 2 & 1 & 0 & 1 & 0 & 14 \\ s_3 & 0 & \boxed{2} & 0 & 0 & 1 & 10 \\ \hline z & -2 & -5 & 0 & 0 & 0 & 0 \end{array}
}

  \item The most negative entry is under $y$, so $y$ enters: a tray of pies pays most. How far can it grow? The ratios of the right-hand side to the positive entries of the column are $(8,\ 14,\ 5)$, the smallest is in the third row, so $s_3$ leaves. Divide that row by the pivot and clear the column.
\handoutdisplay{
\begin{array}{c|ccccc|c}  & x & y & s_1 & s_2 & s_3 & \\ \hline s_1 & \boxed{1} & 0 & 1 & 0 & -0.5 & 3 \\ s_2 & 2 & 0 & 0 & 1 & -0.5 & 9 \\ y & 0 & 1 & 0 & 0 & 0.5 & 5 \\ \hline z & -2 & 0 & 0 & 0 & 2.5 & 25 \end{array}
}

  \item We stand at $(0,\ 5)$ with $z=25$. Only $x$ still has a negative entry. The ratios are $(3,\ 4.5)$, so $s_1$ leaves.
\handoutdisplay{
\begin{array}{c|ccccc|c}  & x & y & s_1 & s_2 & s_3 & \\ \hline x & 1 & 0 & 1 & 0 & -0.5 & 3 \\ s_2 & 0 & 0 & -2 & 1 & 0.5 & 3 \\ y & 0 & 1 & 0 & 0 & 0.5 & 5 \\ \hline z & 0 & 0 & 2 & 0 & 1.5 & 31 \end{array}
}

  \item No entry of the $z$-row is negative: $z=31-2s_1-\tfrac{3}{2}s_3$, and raising a nonbasic variable can only lower $z$. So $(3,\ 5)$ is optimal with $z^*=31$, and the basic $s_2=3$ says that that much dough is left over. Two pivots were enough; entering $x$ first would have gone the long way round through $(7,\ 0)$ and $(6,\ 2)$.
\end{enumerate}
\endgroup
\begin{answerbox}
$(x^*,y^*)=(3,\ 5)$, $z^*=31$.
\end{answerbox}

\problemhead{B2}{Three variables}{On your own}
\textbf{Statement.} Maximize $5x_1+4x_2+3x_3$ subject to $x_1+x_3\le14$, $x_1+3x_2\le6$, $2x_1+x_2+x_3\le26$ and $x\ge0$. Solve by the simplex method with the same entering rule. Give $z^*$, $x_3^*$ and the final value of the third slack $s_3$.

\textbf{Solution.}
\begingroup\sloppy
\begin{enumerate}[label=\textbf{\arabic*.}, leftmargin=*]
  \item The first tableau, with the pivot in a box. The largest coefficient belongs to $x_1$, and the ratios $(14,\ 6,\ 13)$ send out $s_2$.
\handoutdisplay{
\begin{array}{c|cccccc|c}  & x_1 & x_2 & x_3 & s_1 & s_2 & s_3 & \\ \hline s_1 & 1 & 0 & 1 & 1 & 0 & 0 & 14 \\ s_2 & \boxed{1} & 3 & 0 & 0 & 1 & 0 & 6 \\ s_3 & 2 & 1 & 1 & 0 & 0 & 1 & 26 \\ \hline z & -5 & -4 & -3 & 0 & 0 & 0 & 0 \end{array}
}

  \item After the pivot $z=30$. Now only $x_3$ pays, and the ratios $(8,\ 14)$ send out $s_1$.
\handoutdisplay{
\begin{array}{c|cccccc|c}  & x_1 & x_2 & x_3 & s_1 & s_2 & s_3 & \\ \hline s_1 & 0 & -3 & \boxed{1} & 1 & -1 & 0 & 8 \\ x_1 & 1 & 3 & 0 & 0 & 1 & 0 & 6 \\ s_3 & 0 & -5 & 1 & 0 & -2 & 1 & 14 \\ \hline z & 0 & 11 & -3 & 0 & 5 & 0 & 30 \end{array}
}

  \item Every entry of the $z$-row is now nonnegative, so we stop: $x^*=(6,\ 0,\ 8)$, $s_3=6$, $z^*=54$.
\handoutdisplay{
\begin{array}{c|cccccc|c}  & x_1 & x_2 & x_3 & s_1 & s_2 & s_3 & \\ \hline x_3 & 0 & -3 & 1 & 1 & -1 & 0 & 8 \\ x_1 & 1 & 3 & 0 & 0 & 1 & 0 & 6 \\ s_3 & 0 & -2 & 0 & -1 & -1 & 1 & 6 \\ \hline z & 0 & 2 & 0 & 3 & 2 & 0 & 54 \end{array}
}

\end{enumerate}
\endgroup
\begin{answerbox}
$x^*=(6,\ 0,\ 8)$, $z^*=54$, $s_3=6$.
\end{answerbox}

The method needs a first corner. For $Ax\le b$ with $b\ge0$ the origin works. Otherwise make the right-hand sides nonnegative by changing the sign of rows, add an artificial variable to each row that lacks a basic variable and, in Phase I, minimize their sum $w$ with the same pivots. If the minimum is positive the LP is infeasible. If it is zero the point is a corner of the original problem: an artificial variable still basic at value zero is pivoted out, or its row is dropped as redundant, and Phase II starts from there with the real objective.

\problemhead{B3}{A start that has to be found}{On your own}
\textbf{Statement.} Maximize $2x+y$ subject to $x+y\le6$, $x-y\ge2$ and $x,y\ge0$. The origin is infeasible, so run Phase I with one artificial variable $a$ in the second row, then Phase II with the same entering rule. Give the value of $x$ at the end of Phase I, the optimal value $z^*$ and the final value of the surplus $s_2$.

\textbf{Solution.}
\begingroup\sloppy
\begin{enumerate}[label=\textbf{\arabic*.}, leftmargin=*]
  \item With the rows $x+y+s_1=6$ and $x-y-s_2+a=2$ the basis is $s_1,a$. Minimizing $w=a$ means maximizing $-w$, and removing the basic $a$ from that row gives
\handoutdisplay{
\begin{array}{c|ccccc|c}  & x & y & s_1 & s_2 & a & \\ \hline s_1 & 1 & 1 & 1 & 0 & 0 & 6 \\ a & \boxed{1} & -1 & 0 & -1 & 1 & 2 \\ \hline -w & -1 & 1 & 0 & 1 & 0 & -2 \end{array}
}

  \item One pivot in the column of $x$ removes $a$ and gives $w=0$, so the problem is feasible and Phase I ends at the point $(2,\ 0)$.
\handoutdisplay{
\begin{array}{c|ccccc|c}  & x & y & s_1 & s_2 & a & \\ \hline s_1 & 0 & 2 & 1 & 1 & -1 & 4 \\ x & 1 & -1 & 0 & -1 & 1 & 2 \\ \hline -w & 0 & 0 & 0 & 0 & 1 & 0 \end{array}
}

  \item Delete the column of $a$, restore the real objective and price out the basic $x$: $z=4+3y+2s_2$.
\handoutdisplay{
\begin{array}{c|cccc|c}  & x & y & s_1 & s_2 & \\ \hline s_1 & 0 & \boxed{2} & 1 & 1 & 4 \\ x & 1 & -1 & 0 & -1 & 2 \\ \hline z & 0 & -3 & 0 & -2 & 4 \end{array}
}
 Then $y$ enters and $s_1$ leaves; the only positive entry is in its row.
  \item Now $z=10$ at $(4,\ 2)$, and the $z$-row still asks for $s_2$. Its ratio test keeps only the first row, so $y$ leaves again.
\handoutdisplay{
\begin{array}{c|cccc|c}  & x & y & s_1 & s_2 & \\ \hline y & 0 & 1 & 0.5 & \boxed{0.5} & 2 \\ x & 1 & 0 & 0.5 & -0.5 & 4 \\ \hline z & 0 & 0 & 1.5 & -0.5 & 10 \end{array}
}

  \item Done:
\handoutdisplay{
\begin{array}{c|cccc|c}  & x & y & s_1 & s_2 & \\ \hline s_2 & 0 & 2 & 1 & 1 & 4 \\ x & 1 & 1 & 1 & 0 & 6 \\ \hline z & 0 & 1 & 2 & 0 & 12 \end{array}
}
 The optimum is $(6,\ 0)$ with $z^*=12$ and $s_2=4$. A variable that entered may leave later; nothing in the method forbids it.
\end{enumerate}
\endgroup
\begin{answerbox}
Phase I ends at $(2,\ 0)$; $x^*=(6,\ 0)$, $z^*=12$, $s_2=4$.
\end{answerbox}

\begin{trapbox}
In the ratio test negative and zero entries are skipped. Dividing by a negative entry gives a smaller number, and taking it makes the point infeasible.
\end{trapbox}

\begin{selfcheckbox}
In B1 the $z$-row of the last tableau has zeros under $x$, $y$ and $s_2$. Does that mean the optimum is not unique?
\selfcheckanswer No. Those variables are basic, and basic variables always have zeros there. Only a zero under a nonbasic variable can point to another optimum.
\end{selfcheckbox}

% =====================================================================
\section{Duality}\label{sec:duality}

Where does an upper bound on the optimum come from? Take nonnegative weights $y$ and add up the rows of $Ax\le b$. If the weighted row dominates the objective, $A^{\top}y\ge c$, then $c^{\top}x\le y^{\top}Ax\le b^{\top}y$ for all feasible $x$. Looking for the best such bound is an LP in itself.

\begin{definitionbox}{Definition: Dual problem.}
The dual of $\max c^{\top}x$, $Ax\le b$, $x\ge0$ is $\min b^{\top}y$, $A^{\top}y\ge c$, $y\ge0$. A row of the primal gives a variable of the dual, a column gives a row. An equality row gives a free dual variable and a free primal variable gives an equality row. The dual of the dual is the primal. This is the Lagrangian dual of seminar 8 written out for a linear objective.
\end{definitionbox}

\begin{theorembox}{Theorem: Weak and strong duality.}
For feasible $x$ and $y$ we have $c^{\top}x\le b^{\top}y$. If one of the two problems has an optimal solution, so does the other, and the optimal values are equal. If one is unbounded, the other is infeasible.
\end{theorembox}

\begin{proofbox}
Weak duality is the chain above. For the strong part, run the simplex method to an optimal tableau with basis $B$ and put $y^{\top}=c_B^{\top}B^{-1}$. The entries of the $z$-row under the slacks are $y$ and those under the $x_j$ are $(A^{\top}y-c)_j$; optimality says that both are nonnegative, so $y$ is dual feasible. Moreover $b^{\top}y=c_B^{\top}B^{-1}b=c^{\top}x$.
\end{proofbox}

\begin{theorembox}{Corollary: Complementary slackness.}
Feasible $x$ and $y$ are both optimal exactly when $y_i(b-Ax)_i=0$ for every row $i$ and $x_j(A^{\top}y-c)_j=0$ for every column $j$.
\end{theorembox}

\begin{proofbox}
The gap is $b^{\top}y-c^{\top}x=y^{\top}(b-Ax)+x^{\top}(A^{\top}y-c)$, a sum of nonnegative products. It vanishes exactly when every product does.
\end{proofbox}

\begin{whybox}
These are the KKT conditions of an LP: primal feasibility, dual feasibility and complementary slackness, with stationarity built into the dual rows. For an LP no regularity condition is needed.
\end{whybox}

\problemhead{C1}{What an hour of oven time is worth}{On the board}
\textbf{Statement.} Take the bakery of B1 with its optimum $(3,\ 5)$. Write the dual, find $y^*$ from complementary slackness and compare it with the last tableau. For which values of the oven limit $b_1$ does the basis stay optimal? Give $y_1^*$, $y_3^*$ and the largest such $b_1$.

\textbf{Solution.}
\begingroup\sloppy
\begin{enumerate}[label=\textbf{\arabic*.}, leftmargin=*]
  \item One price per resource, one row per product: minimize $8y_1+14y_2+10y_3$ subject to $y_1+2y_2\ge2$, $y_1+y_2+2y_3\ge5$, $y\ge0$. Read it as a buyer who rents all resources: each product must be paid at least what it would earn.
  \item Complementary slackness at $(3,\ 5)$. The dough row is slack, $s_2=3$, so $y_2=0$. Both $x$ and $y$ are positive, so both dual rows are tight: $y_1=2$, then $y_1+2y_3=5$ gives $y_3=1.5$. So $y^*=(2,\ 0,\ 1.5)$, the entries under $s_1,s_2,s_3$ in the $z$-row of the last tableau.
  \item The dual value is $b^{\top}y^*=31=z^*$, so both solutions are certified optimal at once. An extra hour of oven is worth $2$, extra dough nothing, a unit of filling $1.5$.
  \item Change $b_1$ to $8+\delta$. The column of $s_1$ in the last tableau gives how the basic variables move: $x=3+\delta$, $s_2=3-2\delta$, $y=5$. They stay nonnegative for $-3\le\delta\le1.5$, that is for $b_1$ in $[5,9.5]$. In this range $p^*$ has slope $y_1^*$; past its ends the basis changes and so does the price.
\end{enumerate}
\endgroup
\begin{answerbox}
$y^*=(2,\ 0,\ 1.5)$, the basis stays optimal for $b_1\in[5,9.5]$.
\end{answerbox}

\problemhead{C2}{Solve the dual instead}{On your own}
\textbf{Statement.} Maximize $9x_1+10x_2+3x_3$ subject to $2x_1+2x_2+5x_3\le14$, $3x_1+4x_2+2x_3\le24$ and $x\ge0$. The dual has two variables: solve it graphically, then recover the primal solution from complementary slackness. Give $y_1^*$, the common optimal value and $x_1^*$.

\textbf{Solution.}
\begingroup\sloppy
\begin{enumerate}[label=\textbf{\arabic*.}, leftmargin=*]
  \item The dual: minimize $14y_1+24y_2$ subject to $2y_1+3y_2\ge9$, $2y_1+4y_2\ge10$, $5y_1+2y_2\ge3$, $y\ge0$.
  \item Its corners are $(0,\ 3)$, $(3,\ 1)$ and $(5,\ 0)$, with values $(72,\ 66,\ 70)$. The smallest is $66$ at $y^*=(3,\ 1)$. The third row is not needed there: it never binds.
  \item At $y^*$ the third dual row is slack, so complementary slackness gives $x_3=0$; both $y_i>0$, so both primal rows are tight. From $2x_1+2x_2=14$ and $3x_1+4x_2=24$ we get $x^*=(4,\ 3,\ 0)$, and the primal value equals $66$ as it must.
\end{enumerate}
\endgroup
\begin{answerbox}
$y^*=(3,\ 1)$, $x^*=(4,\ 3,\ 0)$, optimal value $66$.
\end{answerbox}

% =====================================================================
\section{Shadow prices and ranges}\label{sec:sensitivity}

Keep the optimal basis and change $b$ a little. The basic values $x_B=B^{-1}b$ move linearly, and while they stay nonnegative the basis stays optimal: the $z$-row does not depend on $b$. So the optimal value is $y^{\top}b$ on that range, with the same $y$.

\begin{theorembox}{Theorem: Shadow prices.}
For a nondegenerate optimal basis, $\partial p^*/\partial b_i=y_i^*$. Globally $p^*(b)$ is piecewise linear and concave; each piece belongs to one basis, and at a kink the left and right slopes are two different optimal prices.
\end{theorembox}

\begin{itemize}
  \item Range of $b_i$: the column of $s_i$ in the final tableau is the $i$-th column of $B^{-1}$, so $x_B+\delta B^{-1}e_i\ge0$ gives the interval for $\delta$.
  \item Range of $c_j$: the point stays and the $z$-row changes. Recompute the row with the new coefficient and require that all its entries stay nonnegative.
  \item Inside a range of $b$ the prices stay and the plan moves; inside a range of $c$ the plan stays and the prices move.
\end{itemize}

\handoutfigure{0.72}{fig_c1}{The optimal value of B1 against $b_1$. The slope is the shadow price of oven time, and it falls at each kink.}

\problemhead{C3}{How far the plan survives}{On your own}
\textbf{Statement.} Maximize $4x+3y$ subject to $x+y\le9$, $2x+y\le14$, $y\le6$ and $x,y\ge0$. Find the optimum and the prices. Then give $y_1^*$, the largest $b_1$ (the right-hand side of the first row) for which the same basis stays optimal, and the largest coefficient of $x$ in the objective for which the same point stays optimal.

\textbf{Solution.}
\begingroup\sloppy
\begin{enumerate}[label=\textbf{\arabic*.}, leftmargin=*]
  \item The first two rows meet at $x^*=(5,\ 4)$, the value is $32$ and the third row has slack $2$. The prices solve $y_1+2y_2=4$, $y_1+y_2=3$: $y^*=(2,\ 1,\ 0)$, with zero for the third row.
  \item With $b_1=9+\delta$ the two tight rows give $x=5-\delta$ and $y=4+2\delta$. Nonnegativity alone would allow $\delta\le5$, but the third row asks for $y\le6$, that is $\delta\le1$. So the basis holds up to $b_1=10$, and past that point the row $y\le6$ takes over.
  \item Replace the coefficient of $x$ by $c_1$. The point stays optimal while the level line lies between the two tight rows: their slopes require $1\le c_1/3\le2$, so $c_1\in[3,6]$. The same follows from $y_1=2\cdot3-c_1\ge0$ and $y_2=c_1-3\ge0$.
\end{enumerate}
\endgroup
\begin{answerbox}
$y_1^*=2$, the basis holds up to $b_1=10$, the point up to a coefficient $6$.
\end{answerbox}

\begin{trapbox}
A shadow price is a derivative, not a market price. It says nothing beyond its range, and at a degenerate corner the prices need not be unique, so the number may depend on the solver.
\end{trapbox}

\begin{selfcheckbox}
In C1 the price of oven time is $2$ on $[5,9.5]$. What are the slopes of $p^*$ below and above this range?
\selfcheckanswer Below it $5$, because the oven then limits pies alone; above it $0$, because the oven is no longer short.
\end{selfcheckbox}

% =====================================================================
\section{Degeneracy, cycling and running time}\label{sec:failures}

The finite step count of the simplex method rests on one fact: every pivot strictly raises $z$, so no basis repeats. That fails when a basic variable is zero, because the ratio test may then give a step of length zero. Such a degenerate pivot changes the basis without moving the point.

\begin{theorembox}{Theorem: Bland's rule.}
If among all columns with a negative entry the entering variable has the smallest index, and among the rows with the smallest ratio the leaving variable has the smallest index, the simplex method stops after finitely many pivots.
\end{theorembox}

The rule is slow in practice, so codes use it as a fallback or perturb $b$ slightly to break the ties. As for speed, the Klee–Minty cube shows that Dantzig's rule can visit all $2^n$ corners, and no pivot rule is known to be polynomial. Yet on real problems the number of pivots is usually a small multiple of $m$. Polynomial algorithms for LP exist: the ellipsoid method and interior-point methods, and the latter compete with the simplex method in practice.

\handoutfigure{0.72}{fig_km}{Pivots of Dantzig's rule on Klee–Minty cubes and on random problems of the same sizes.}

\begin{itemize}
  \item \texttt{linprog} status 3: unbounded. Look for a forgotten constraint.
  \item \texttt{linprog} status 2: infeasible. Look for rows that contradict each other; the certificate of Farkas' lemma names them.
  \item \texttt{res.ineqlin.marginals} holds the prices of a minimization: for a problem turned around by a sign they change sign too.
\end{itemize}

% =====================================================================
\section{Why integer optimization needs all this}\label{sec:integer}

A solver for integer models spends most of its time solving LPs. The relaxation of seminar 8 is an LP, and its value is the bound against which every integer point is judged. Strong duality makes that bound cheap to certify: any dual feasible $y$ gives it without trusting the solver.

The tableau helps as well. A new bound on a variable or a new valid row leaves the old basis dual feasible, so the next LP restarts from there, usually after a few pivots, instead of from scratch. And when the optimal corner of a relaxation happens to be integer, nothing more needs to be done: the search is over for that subproblem.

\end{document}
